\documentclass[10pt]{article}
\usepackage[utf8]{inputenc}
\usepackage[T1]{fontenc}
\usepackage{amsmath}
\usepackage{amsfonts}
\usepackage{amssymb}
\usepackage[version=4]{mhchem}
\usepackage{stmaryrd}
\usepackage{bbold}
\usepackage{graphicx}
\usepackage[export]{adjustbox}
\graphicspath{ {./images/} }

\DeclareUnicodeCharacter{21D2}{\ifmmode\Rightarrow\else{$\Rightarrow$}\fi}
\DeclareUnicodeCharacter{2192}{\ifmmode\rightarrow\else{$\rightarrow$}\fi}

\begin{document}
Competitive equilibrium - intro\\
The optimal growth problem studied so far is the problem of the Social Planner\\
(who allocates resourses to max agents' utility s.t. the feasibility constraint of the economy)

The Social Planner das not rely on markets lall the markets are shut down, social plasmer gust tells everybody what to to).

The social Planner carnot exist in reality.\\
Nevertheless, finding the social Planner's allocation is very useful - because it describes the most efficient allocation of resources in the economy.\\
( = Pareto Optimal)\\
(= "the efficient)\\
In practice, the economic agents interact through the markets ⇒\\
⇒ we want to be able to

\begin{enumerate}
  \item model these interactions
  \item compare the market butcomes woth efficient allocations
  \item analyze policies / and design the ones that improve efficiency)
  \item quantify the losses associated with economic frictions and design the policies mitigating them.
\end{enumerate}

The general structure of the equilibrium (not rigorous):\\
Consumers: max life-time utility\\
choose. how much to work $\Rightarrow L^{5}$ (labor supply).

\begin{itemize}
  \item how much to save $\Rightarrow K^{\prime}$ (capital supply)
  \item how much to consume $\Rightarrow C^{d}$ (cons. demand)
\end{itemize}

Firms: max life-time profit\\
choose a inputs: how much to hire $\Rightarrow 2^{d}$ (labor demand)\\
how much capital to rent $\Rightarrow K^{d}$\\
ontport: how much to produce $\Rightarrow Y^{S}$

In the equilibrium, the price are such that all markets clear In general, any competitive equilibrium is defined as a pair (vector) ( $p, q$ ): prices and quantities, such that

\begin{enumerate}
  \item Given $p$, households chasse $q$ optimally "competitive"
  \item Given $p$, firms choose $q$ aptimally
  \item Prices $p$ clear the markets ( $q$ is the same in 1) 621)
\end{enumerate}

Since we may work isth two setups, sequential and recursive, we introduce two types of equilitia,

\begin{itemize}
  \item sequential competitive equilibrium and
  \item recursive competitive equilibrium
\end{itemize}

We will start with\\
→ Sequential competitive equilibrium there are two approaches: (diftering in time of trade):

\begin{enumerate}
  \item Arrow-Debren time-o trading:
\end{enumerate}

\begin{itemize}
  \item all the goods, for all the periods, are traded in period o. from $t=1$ onwards, agents pust take actions/make exchanges corresponding to promices made at $t=0$
\end{itemize}

\begin{enumerate}
  \setcounter{enumi}{1}
  \item Sequential trading
\end{enumerate}

\begin{itemize}
  \item markits reopen in every period
\end{itemize}

The structure of the economy:

\section*{Consumers:}
\begin{itemize}
  \item own capital → rent it to the firms
  \item have 1 unit of time endowment $\left(n_{+}+l_{+} \leq 1\right)$\\
→ labor → brings income / supply to leisure → brings willity the firms)
  \item buy consumption grodds from the market (ca) (derive utility from consumption)
  \item hold ownership shares in the firms (representative portfolio)
\end{itemize}

\section*{→ receive dividends}
\begin{itemize}
  \item consumption, capital, tividends are perfect substitutes.
\end{itemize}

\section*{God: to maximize life-time utility}
\section*{Frus:}
\begin{itemize}
  \item owned by consumers → distribute all the profits as dividends\\
-choose inputs: capital to rent labor to hire
\end{itemize}

$$
\} \Rightarrow F\left(k_{+}, n_{+}\right)
$$

$$
\text { ⇒ profit }=\text { output }- \text { all costs }
$$

Goal: maximize value (total profit)\\
Markets: there are 3 types of goods in this economy

\begin{itemize}
  \item "Promisis" of supplying final good $y_{t}$\\
price $=p_{t}, t \geqslant 1$\\
normalize $\rho_{0}=0$
  \item "Promises" of providing labor services $x_{1}$ :\\
price $=\omega_{+}$in terms of good-t units\\
$\Rightarrow p_{+} \omega_{+} \cdot n_{+}$paid to the consumer in period. if helshe promices to work $n_{+}$in period $t$
  \item "Promises" of renting capital $k_{+}$:\\
price $=r_{t}$, in terms of period-t goods\\
$\Rightarrow p_{+} \cdot r_{+} \cdot k_{+}$paid to the cons. at $t=0$ if he/she promises to rent $k_{+}$units of capital in periodt\\
(A) Consumer's problem:
\end{itemize}

$$
\begin{aligned}
\left\{\max _{\left\{k_{+1}, c_{+}, n_{+}, 4_{3}\right\}_{+=0}^{\infty}}\right. & \sum_{+=0}^{\infty} \beta^{+} u\left(c_{1}, l_{\infty}\right) \\
\text { s.t. } & l_{+}+n_{+} \leq 1, \quad l_{+} \geq 0, n_{+} \geq 0 \\
& k_{+1} \geq 0, c_{+} \geq 0 \\
& k_{0}-\text { given } \\
& \text { AND date-0 budget constraint }
\end{aligned}
$$

the cost of purchasing all promises att-o $\leq$ the revenue from selling all promises

\section*{Revenues come from:}
\begin{itemize}
  \item selling promises of future labor senices: $\sum_{+=0}^{\infty} p_{+} w_{+} n_{+}$
  \item selling promises to rent cajorital: $\sum_{+=0}^{\infty} p_{+} r_{+} k+$
  \item selling promises to deliver final goods to other consumers:
\end{itemize}

$$
\sum_{t=0}^{\infty} p_{t}(1-\infty) k_{t}
$$

\begin{itemize}
  \item revenues from form ownership
\end{itemize}

$$
\Pi=\sum_{t=1}^{\infty} p+\pi
$$

Resourses spent on

\begin{itemize}
  \item buying promices to deliver goods (cons. and capoital
\end{itemize}

$$
\sum_{t=0}^{\infty} p_{+} G_{+}+\sum_{t=0}^{\infty} p_{+} k_{++1}
$$

⇒ Date-o budget constraint is

$$
\sum_{t=0}^{\infty} p_{+}\left(G_{+}+k_{++1}\right) \leqslant \sum_{t=0}^{\infty} p_{+}\left(r_{+} k_{+}+\omega_{+} n_{+}+(1-\alpha) k_{+}+\pi_{+}\right)
$$

Given $\left\{p_{+}, r_{+}, w_{+}\right\}_{+=0}^{\infty}$ and $\left\{\pi_{+}\right\}_{+=0}^{\infty}$, the cons.\\
choose $\left\{G_{1}, n_{+}, l_{+}, k_{+k+1}\right\}_{+=0}^{\infty}$\\
(B) Firm's problem:

Goal: to maximize the "value" of the firm = profit quevated by sellny/buyn'p doums at to 0\\
$\operatorname{Max}_{\left\{k_{+}, n_{+}\right\}_{+=0}^{\infty}} \sum_{+=0}^{\infty} p_{+}\left[F\left(k_{+}, n_{+}\right)-r_{+} k_{+}-\omega_{+} n_{+}\right]$\\
$p_{+} F\left(k_{r}, n_{+}\right)$- revenues from selling promises to deliver goods in periodt\\
$P_{+} r_{-} k_{+}$-cost of buying promises to rent cospiral\\
$p_{+} \omega_{T} n_{+}$- cost of buying promises to use labor\\
The firm's decision problem can be decomposed into a set of static problems:\\
$\pi_{+}=\max _{k_{+}, n_{+}} F\left(k_{+}, n_{+}\right)-r_{+} k_{+}-\omega_{+} n_{+}, \quad t \geqslant 0$\\
$\Rightarrow M=\sum_{+=0}^{\infty} p+\pi_{+}$no is constraints. to shareholders in period o

⇒ Given $\left\{r_{+}, \omega_{+}\right\}_{+=0}^{\infty}$, we can denve $\left\{k_{+}, h_{+}, \pi_{+}\right\}_{+=0}^{\infty}$\\
(c) A Competitive Equilibrius in (formal definition) consists of prices $\left\{p_{+}^{*}, r_{+}^{*}, w_{+}^{*}\right\}_{+=0}^{\infty}$ and quantities $\left(G_{+}^{*}, k_{+}^{*}, n_{+}^{*}, \pi_{+}^{*}\right\}_{+=0}^{\infty}$ such that

\begin{enumerate}
  \item Given prices $\left\{p_{+}^{*}, r_{+}^{*}, \omega_{+}^{*}\right\}$ and ofividends $\left\{\pi_{+}^{*}\right\}$, the quantities $\left\{G^{*}, k_{*}^{*}, n_{*}^{*}\right\}$ solve the cons'dicision problem
  \item Given prices $\left\{p_{+}^{*}, r_{+}^{*}, \omega_{+}^{*}\right\}$, the quantities $\left\{k_{*}^{*}, n_{*}^{*}, \pi_{t}^{*}\right\}$ solve the firm's decision problem\\
$\left\{k_{*}^{*}, n_{*}^{*}, \pi_{t}^{*}\right\}$ 'solve the firm's decision problem
  \item Prices are such that all markets clear:
\end{enumerate}

\begin{itemize}
  \item markets for copital cend labor promises clear sirce the same $n_{+}^{*}, k_{+}^{*}$ solve cons' and from' problems
  \item the market for froal good promises must clear
\end{itemize}

$$
\underset{\substack{\text { promises baight } \\ \text { by consumers }}}{L_{+}^{*}+k_{++1}^{*}}=\underbrace{F\left(k_{+}^{*}, n_{+}^{*}\right)}_{\substack{\text { promises sold } \\ \text { by firms }}}+\underbrace{(1-\sigma) k_{+}^{*}}_{\substack{\text { proming } \\ \text { sold by consumers }}}, t \geq 0
$$

Our objectives are: to solve for the competritive equil. (CE) compare the CE allocation with the efficient allocation

Let's consider the simplified model, which is similar to the Optimal Growth model considered so far: $u\left(c_{1} l_{+}\right)=u\left(c_{+}\right) \quad$ (no disutitity from work $\Rightarrow n_{+}=1, l_{+}=0$ )\\
$\left.(C P): \max _{\left\{k_{+1}, c_{+}\right\}} \sum_{+=0}^{\infty} \beta^{+} u / c_{+}\right)$

$$
\begin{array}{lll}
\text { s.t. } & \sum_{t=0}^{\infty} p_{t}\left(c+k_{+\cdots}\right) \leqslant \sum_{t=0}^{\infty} p_{t}\left(r_{+} k_{+}+(1-\delta) k_{+}+\omega_{+}+\pi_{+}\right) \\
& c_{+} \geq 0, t \geqslant 0 & \left(\eta_{t}\right) \\
& k_{+1} \geqslant 0, t \geqslant 0 & \left(\nu_{t}\right) \\
& k_{0} \text {-given } &
\end{array}
$$

$$
\begin{aligned}
L=\sum_{t=0}^{\infty} \beta^{t} u\left(c_{t}\right) & -\lambda \sum_{t=0}^{\infty} p_{t}\left(c_{t}+k_{t+1}-r_{t} k_{t}-(1-\sigma) k_{t}-\omega_{t}-\pi_{t}\right) \\
& +\sum_{t=0}^{\infty} \eta_{t} c_{t}+\sum_{t=0}^{\infty} \nu_{t} \cdot k_{t+1}
\end{aligned}
$$

$K-T$ conditions:

$$
\begin{array}{lll}
\text { FOC: } & \left(c_{+}\right): \beta^{+} u^{\prime}\left(c_{+}\right)-\lambda p_{+}+\eta_{+}=0 & , t \geqslant 0 \\
& \left(k_{++1}\right):-\lambda \cdot p_{+}+\lambda p_{++1}\left(1+r_{++1}-\alpha\right)+\nu_{+}=0, & t \geqslant 0 \\
K-T: & \lambda \geqslant 0 & \lambda \cdot \sum_{+=0}^{\infty} p_{+}\left(c_{+}+k_{++1}-\left(1+r_{+}-\alpha\right) k_{+}-\omega_{+}-\pi_{+}\right)=0 \\
& \eta_{+} \geqslant 0 & \eta_{+} \cdot c_{+}=0 \\
& \nu_{+} \geqslant 0 & \nu_{+} \cdot k_{+++}=0
\end{array}
$$

Can we eliminate corners $c_{4}=0$ and $k_{41}=0$ ?

\begin{enumerate}
  \item If $u^{\prime}(0)=+\infty \Rightarrow$\\
$C_{+}=0$ is impossible since $\lambda$ and $\eta_{+}$are scolars $\Rightarrow \eta_{t}=0 \Rightarrow \lambda>0$
  \item If can argue that $k_{++1}=0 \Rightarrow r_{++1}=\infty$
\end{enumerate}

FOC ( $k_{++1}$ ): $p_{+}=p_{+1}\left(1+r_{++1}-\alpha\right)+\frac{\nu_{t}}{\lambda} \geqslant p_{++1}\left(1+r_{+1}-\alpha\right)$\\
When it holds with " $=$ ": $p_{+}=p_{++1}\left(1+r_{++1}-\delta\right)$ no-arbitrage condition?\\
not possible to gain by simultamedidy

\begin{itemize}
  \item buying promises to deliver a unit of final good in period $t$
  \item selling promizes to rent a unit of capartal in period $t+1$\\
If " $>$ " $\Rightarrow$ want to sell promises to rent $k_{x+1} \Rightarrow \Rightarrow k_{a n} \geqslant 0$ binds.\\
If $r_{1+1}=\infty$ when $k_{a+1}=0 \Rightarrow$ this FOC cannot be\\
satisfied for $\min _{t}$ for which $k_{t+1}=0 \quad($ skalar $\geq \infty)$\\
$\Rightarrow \nu_{+}=0$
\end{itemize}

$$
\Rightarrow|F O C| \Rightarrow \begin{cases}\beta^{+} u^{\prime}\left(c_{4}\right)=\lambda p_{+} & \Rightarrow \frac{u^{\prime}\left(c_{+}\right)}{\beta u^{\prime}\left(c_{++1}\right)}=\frac{p_{+}}{p_{+1}} \\ & (\forall t \geq 0)\end{cases}
$$

Euler Equation (like in the OGM with linear return)\\
The consumer observes $\left\{r_{+}\right\}_{t \geq 0}$ and chooses $\{G\}_{t \geq 0}$\\
⇒ EE is the first order difference equation w.r.t. \{Ct\}.\\
If $C_{0}$ is known $\Rightarrow$ can get the whole signence, $C_{t}\left(C_{0}\right)$\\
How to fird Co? Use the budget constraint!

$$
\begin{gathered}
\sum_{t=0}^{\infty} p_{+} G_{+}\left(C_{0}\right)+\sum_{+=0}^{\infty} p_{+} k_{++1}=\sum_{+=0}^{\infty} \underbrace{p_{+}\left(1+r_{+}-\infty\right)}_{=p_{+-1}, t \geq 1} k_{+}+\sum_{+=0}^{\infty} p_{+}\left(\omega_{+}+\pi_{+}\right) \\
\text {come from } E E \\
p_{0}=1 \text { (normalize) }
\end{gathered}
$$

$$
\begin{aligned}
F O C \Rightarrow p_{+}\left(1+r_{+}-\alpha\right) & =p_{+-1} \\
\Rightarrow \sum_{+=0}^{\infty} p_{+} c_{+}\left(c_{0}\right) & =\lim _{T \rightarrow \infty} \underbrace{\sum_{+=0}^{T} p_{+}\left(1+r_{+}-\alpha\right) k_{+}}_{p_{0}\left(1+r_{0}-\alpha\right) k_{0}+\sum_{+=0}^{T-1} p_{+} k_{++1}}-\sum_{+=0}^{T} p_{+} k_{++1})+\sum_{+=0}^{\infty} p_{+}\left(\omega_{+}+\pi_{+}\right) \\
& =\underbrace{p_{0}\left(1+r_{0}-\alpha\right) k_{0}}_{\text {given to the cons. }}-\lim _{T \rightarrow \infty} p_{T} k_{T+1}+\underbrace{\sum_{+=0}^{\infty} p_{+}\left(\omega_{+}+\pi_{+}\right)}_{\text {given to the }}
\end{aligned}
$$

The consumer chooses $\left\{k_{t+1}, C_{4}\right\}_{+\geq 0}$ to max. life-time utility $\left\{c_{C}\right\}$ is pinned down by $c_{0}$, increasing in $c_{0}$ ⇒ chose $\left\{k_{t+1}\right\}_{t \geqslant 0}$ such that $\lim _{t \rightarrow \infty} p_{t} k_{t+1}=0$

$$
\left(p_{+} \geq 0, k_{++1} \geq 0\right)
$$

Thus, conditional on prices satisfying $p_{+}^{*}\left(1+r_{+}^{*}-d\right)=p_{+-1}^{*},-1 \geqslant 1$ the solution to the consumer's problem is: I (NA) $\left\{c_{+}^{*}, k_{++1}^{*}\right\}_{+=0}^{\infty}$ such that

$$
\begin{aligned}
& u^{\prime}\left(c^{*}\right)=\beta\left(1+r_{++1}^{*}-\delta\right) u^{\prime}\left(c_{+1}^{*}\right) \\
& \sum_{t=0}^{\infty} p_{+} c_{+}^{*}=p_{0}^{*}\left(1+r_{0}^{*}-\delta\right) k_{0}+\sum_{t=0}^{\infty} p_{+}^{*}\left(\omega_{+}^{*}+\pi_{+}^{*}\right) \\
& \lim _{t \rightarrow \infty} p_{+}^{*} k_{++1}^{*}=0
\end{aligned}
$$

(FP) $\max _{k_{+}, n_{+}} \sum_{+\infty}^{\infty}\left(F\left(k_{+}, n_{+}\right)-\omega_{+} n_{+}-r_{+} k_{+}\right), \forall t \geq 0$

FOC:

\[
\begin{array}{ll}
\left(k_{+}\right) & F_{1}\left(k_{+}^{*}, n_{+}^{*}\right)=r_{+}^{*}  \tag{2}\\
\left(n_{t}\right) & F_{2}\left(k_{+}^{*}, n_{+}^{*}\right)=\omega_{+}^{*}
\end{array}
\]

$$
\Rightarrow \pi_{+}^{*}=F\left(k_{+}^{*}, n_{+}^{*}\right)-r_{+}^{*} k_{+}^{*}-\omega_{+}^{*} n_{+}^{*}
$$

\section*{Assumptions on $F(,, i)$ :}
\begin{enumerate}
  \item $F_{i}^{\prime}\left(R_{+}, R_{+}\right) \rightarrow \mathbb{R}_{+}$
  \item $F$-monotone in each argument
  \item $F$-strictly concave in each argument
  \item $F_{1}(0, n)=+\infty \quad \Rightarrow \quad k_{++1}^{*}=0$ is impossible
\end{enumerate}

$$
\left(\Rightarrow r_{++1}^{*}=\infty\right)
$$

\begin{enumerate}
  \setcounter{enumi}{4}
  \item If $u(a=l)$ depends on leisure ⇒
\end{enumerate}

⇒ for the same reason assume $F_{2}(k, 0)=+\infty$

The competitive eopiritibrium is

$$
\left\{p_{+}^{*}, r_{+}^{*}, w_{+}^{*}\right\}_{0}^{\infty}\left\{\bar{s}_{t}^{*}\right\}_{0}^{\infty},\left\{G^{*}, n_{+}^{*}, k_{++1}^{*}\right\}_{0}^{\infty}:
$$

\begin{itemize}
  \item conditions (1) hold I Cons. behaves upamally green prices)
  \item conditions (c) hold (Frm.-"-)
  \item markets clear:
\end{itemize}

$$
n_{+}^{*}=1, \quad t \geq 0
$$

kiti the same in (9) and (2) (capital mets) $c_{+}^{*}+k_{++}^{*}=F\left(k_{+}^{*}, n_{+}^{*}\right)+(1-\alpha) k_{-}^{*}$ (goodes mk)

In general equilibrium is not easy to find (how to get the sequence of equilibr prices?)

It's much easier to solve the Social planner's problem (very common)\\
⇒ How does CE compare with the SP's allocation?


\begin{align*}
(S P): & u^{\prime}(\underbrace{F\left(k_{+}, 1\right)+(1-\alpha) k_{+}-k_{+}+1}_{C_{+}})=\beta(\underbrace{\left.F_{1}\left(k_{++1}, 1\right)+1-\alpha\right)}_{f^{\prime}\left(k_{++1}\right)} u^{\prime}(\underbrace{\left(F\left(k_{++1}\right)+(1-\alpha) k_{+1}-k_{+2}\right.}_{C_{++1}}) \\
& \lim _{t \rightarrow \infty} \beta^{+} f^{\prime}\left(k_{+}\right) \cdot u^{\prime}\left(G_{1}\right) \cdot k_{+}=0 \quad(T \vee C) \\
& c_{+}+k_{++1}=F\left(k_{+1}\right)+(1-\alpha) k_{+} \quad \text { (FES) : easibility } \tag{1}
\end{align*}


(CE): conditrous (1) $\left.\quad i^{\prime} / G^{*}\right)=\beta\left(1+r_{+}^{*}-\alpha\right) U^{\prime}\left(C_{+1}^{*}\right)$


\begin{align*}
& \left.\sum_{0}^{\infty} p_{+}^{*} c_{+}^{*}=1+r_{0}^{*}-\alpha\right) k_{0}+\sum_{0}^{\infty} p_{+}^{*}\left(\omega_{+}^{*} \Delta \pi_{+}^{*}\right),  \tag{2}\\
& \lim _{t \rightarrow \infty} p_{+}^{*} k_{++1}^{*}=0 \tag{3}
\end{align*}


conditions (2) $\quad F_{1}\left(k_{+}^{*}, 1\right)=r_{+}^{*}\left(u_{1}\right) \quad \pi_{+}^{*}=F\left(k_{2}^{*}, 1\right)-r_{+}^{*} k_{+}^{*}-w_{2}^{*} n_{+}^{*}$


\begin{equation*}
F_{2}^{*}\left(k_{1}^{*}, 1\right)=\omega_{1}^{*}(5) \tag{6}
\end{equation*}


met dears: $n_{+}^{*}=1$ (7)

$$
k_{+1}^{*}+G^{*}=F\left(k_{+}^{*}, 1\right)+(1-\infty) k_{+}^{*}(8)
$$

\section*{Ist Weltare Theorem:}
A date-o competitive equilibrium is Parets Optimal\\
Proof: $(1)+(4)+(8) \Rightarrow$ EE in the (SP) is satistical (8) ⇒ (FES) is satistied\\
(3) ⇒ TVC in the (SP) is satisfied

$$
\begin{aligned}
& \lim _{t \rightarrow \infty} p_{+}^{*} k_{+1}^{*}=0 \\
& p_{t}=p_{t+1}\left(1+r_{+1}-\sigma\right) \quad \quad \quad\left(\text { FOC w.r.t. }\left(k_{+1}\right)\right) \\
& \Rightarrow \lim _{+\rightarrow \infty} p_{++1}^{*}\left(1+r_{++1}^{*}-\alpha\right) k_{+1}^{*}=0 \\
& \lim _{t \rightarrow \infty} p_{+}^{*} \underbrace{\left(1+r_{+}^{*}-\alpha\right)} k_{+}^{*}=0 \\
& \quad=1+F_{1}\left(k_{+}^{*}, 1\right)-\alpha \quad(\text { by }(4)) \\
& \quad=f^{\prime}\left(k_{+}\right)
\end{aligned}
$$

FOC w.r.t. $\left(c_{+}\right)$in the cons. problem

$$
\beta^{*} u\left(c^{*}\right)=\lambda p_{*}^{*} \quad-1 \quad p_{*}^{*} \quad o u^{\prime}\left(c_{*}^{*}\right)
$$

$$
\begin{aligned}
& \beta^{*} u^{\prime}\left(c^{*}\right)=\lambda p_{+}^{*} \\
& \beta^{*-1} u^{\prime}\left(c_{+-1}^{*}\right)=\lambda p_{+-1}^{*} \Rightarrow \frac{p_{+}^{*}}{p_{+-1}^{*}}=\beta \frac{u^{\prime}\left(c^{*}\right)}{u^{\prime}\left(c_{-1}^{*}\right)} \\
& p_{+}^{*}=\frac{p_{+}^{*}}{p_{+-1}^{*}} \cdot \frac{p_{+-1}^{*}}{p_{+-2}^{*}} \cdot \cdots \frac{p_{-1}^{*}}{p_{0}^{*}} \cdot \underbrace{*}_{0}(\text { normalization }) \\
& \Rightarrow p_{+}^{*}=\beta^{*} \frac{u^{\prime}\left(c_{*}^{*}\right)}{u^{\prime}\left(c_{0}^{*}\right)} \\
& \Rightarrow \lim _{t \rightarrow \infty} p_{+} f^{\prime}\left(k_{*}^{*}\right) \cdot k_{*}^{*}= \\
& =\frac{\lim \beta^{*}\left(u^{\prime}\left(c_{*}^{*}\right) \cdot k_{*}^{*}\right.}{u^{\prime}\left(c_{0}\right)}=0 \Rightarrow(T V C) \operatorname{lod}()
\end{aligned}
$$

2nd Weltare Theorem:

\section*{A Pareto Optimal allocation can be achieved in a competitive equilibrium.}
\section*{Proof:}
Suppose $\left\{G, k_{r+1}\right\}_{+=1}^{\infty}$ satisty (EE), (TVC), (FES).\\
Show that $\exists$ prices such that (1)-(8) are satistied for these $\left\{c_{x}, k_{x+1}\right\}$ of prices.\\
Set prices $\stackrel{\left\{n_{+} r+\right\}_{t=0}^{\infty}}{b^{2}}$ and dividends to satisfy (4)-(6).\\
Then (1) also holds.\\
(FES) implies (8). Obviously, (7) holds.\\
Setting $p_{t}=p_{t-1} \cdot \beta \frac{u^{\prime}(c)}{u^{\prime}(c-1)}$ and recalling that $p_{0}=1$\\
pirs to win equilibr. Ep+3. And guarantices that the rest of the conditions eve satistical.

Frually, the BC, in the cars. decision problem is aptained by multiplying the SP's feasibility constraints by $p_{+}$and summing them up (and use that $\pi_{+}=F\left(k_{+}, n_{+}\right)-\omega_{+} n_{+}-r_{+} k_{+}$)

The same technological environment, just different trading arrangements: trade occurs in every period

Commodity space:

\begin{itemize}
  \item Final goods (can be used for cons. or capital stock)
  \item normalize price to 1
  \item Capital rental services
  \item price $r_{+}$(measured in terms of time-t goods)
  \item Labor services
  \item price cot
\end{itemize}

Consumer's problem:

$$
\begin{aligned}
& \operatorname{Max}_{\imath\left(k_{+}, k_{+}, l_{+}\right\}_{+\infty}^{\infty}} \sum_{+=0}^{\infty} \beta^{+} u\left(c_{+}, l_{+}\right) \\
& \text {s.t. } l_{+}+n_{+} \leq 1 \quad l_{+} \geq 0 \quad n_{+} \geq 0 \\
& k_{++1} \geq 0 \quad c_{+} \geq 0 \quad k_{0}-\text { given } \\
& (\oplus) \quad B \text { adget constr. for every t: } \\
& c_{+}+k_{++1} \leq r_{+} k_{+}+\omega_{+} n_{+}+\pi_{+}+(1-\delta) k_{+}
\end{aligned}
$$

To solve it: (FOC) + (TVC) (like the OGM with linear\\
technology)\\
Given $\left\{\omega_{+}, r_{+}, T_{+}\right\}_{+=0}^{\infty}$ : find optimal $\left\{C_{+}, k_{+\cdots}, n_{+}, l_{n}\right\}_{-\infty}^{\infty}=0$\\
Firm's Problem:

$$
\begin{aligned}
\mathbb{P}_{+} & =\operatorname{Max}_{k_{+}, n_{+}}\left[F\left(k_{+}, n_{+}\right)-r_{+} k_{+}-w_{+} n_{+}\right], \forall t \geq 0 \\
& \Rightarrow F_{1}\left(k_{+}, n_{+}\right)=r_{+} \\
& F_{2}\left(k_{+}, n_{+}\right)=w_{+}
\end{aligned}
$$

Given $\left\{w_{+}, r_{+}\right\}_{+=0}^{\infty}$, quantities $\left\{k_{+}, n_{+}, \pi_{+}\right\}_{+=0}^{\infty}$ solve the firm's problem.

\section*{A competitive equilibrium}
consists of prices $\left\{r_{+}, w_{+}\right\}_{t=0}^{\infty}$ \& quantities $\left\{G_{1} k_{+1}, n_{+}, l_{+}, \pi_{+}\right\}_{+=0}^{\infty}$ such that:

\begin{enumerate}
  \item Given prices and dividends, consumers behave optinally
  \item Given prices, firms behave opormally
  \item prices are such that the markets clear
\end{enumerate}

\begin{itemize}
  \item capital \{ks\}: same in cons' \& firms' problems
  \item lasor \{nt\}: $\_\_\_\_$\\
$\_\_\_\_$\\
$\_\_\_\_$\\
-goods: $c_{+}+k_{++1}=F\left(k_{+}, n_{+}\right)+(1-\alpha) k_{+}$
\end{itemize}

\section*{Remark:}
Since $\pi_{+}=F\left(k_{+}, n_{+}\right)-r_{+} k_{-}-\omega_{+} n_{+}$, the goods met clearing conol. is the same as the cons' budget constr. But this is not true, for example, if there are different types of consumers (e.g. only some of them own firms \& receive dividends)

\section*{Weltare properties:}
HW: establish that the $1^{\text {st }}$ \& $2^{\text {nd }}$ Welfare\\
\includegraphics[max width=\textwidth, alt={}, center]{0d4e3e17-8cdd-4b09-a1b8-f3507f55c416-18_78_208_2275_199}\\
theorems hold.


\end{document}